\clearpage\thispagestyle{empty}
\pdfbookmark[0]{Quellen zum Brief vom 17. August 1810}{gaussexcerpt1810aug-sources}
\begin{center}\textbf{Quellen zum Brief vom 17. August 1810}\end{center}
\noindent Diese KI-Übertragung ins heutige Deutsch folgt dem Brief vom 17. August 1810 in \href{https://gdz.sub.uni-goettingen.de/id/PPN236007467}{Gauß’ \textit{Werken}, Bd.~VI} (Göttingen, 1874), S.~318--319. Der Erstdruck erschien als Beitrag XXXI in der Zeitschrift \href{https://archive.org/download/bub_gb_YCI4AAAAMAAJ/page/n4.jpg}{\textit{Monatliche Correspondenz}, Bd.~XXII} vom September 1810, S.~\href{https://archive.org/download/bub_gb_YCI4AAAAMAAJ/page/n292.jpg}{289} und \href{https://archive.org/download/bub_gb_YCI4AAAAMAAJ/page/n293.jpg}{290}. Das Exemplar der University of Michigan wurde von Google digitalisiert und ist im Internet Archive zugänglich.
\par\smallskip\noindent Die Werkausgabe nennt in der Quellenzeile Band XVIII. Auf dem Titelblatt der Zeitschrift steht Band XXII. Es nennt F. von Zach als Herausgeber. Der Erstdruck setzt zu Beginn vier Punkte und vor der Hoffnung auf den Weiterbau einen Schlusspunkt. Er bezeichnet die beiden in Zeit gemessenen Größen mit Strichen und schreibt Freytagswerder. Die Übertragung folgt der Werkausgabe. KI-Übertragung und KI-Satz.
\clearpage
